% Transcription — fonds Alexandre Grothendieck, shelfmark 28, batch 1 (pages 1-10).
% Established from the facsimile archives/batches/28/batch-01.pdf (a mirror of
% https://grothendieck.umontpellier.fr/28.pdf), per the skill
% transcribe-grothendieck. The folder is ten pages and this is its only batch.
%
% All ten leaves are in his hand, in ink, and all ten are mathematics; none is
% skipped. They fall into two runs, each opened by a heading of his own giving
% a name and a date — « Demazure 5.1.1970 » (pages 1-5) and « Demazure
% 12.2.70 » (pages 6-10). Those two lines are the only dating the leaves carry
% and they are transcribed as the section titles; nothing on the pages says
% what the occasion was, and this file does not say either.
%
% The inventory's title is bracketed, hence the archivists', and it names
% SGA 3. The pages themselves keep to one question: the algebraic subgroups of
% the Cremona group. The first run sets up pseudo-morphisms and the group
% functors Psaut and Baut, with Cr_{n,k} = Baut_k(k[t_1,...,t_n]), then
% pseudo-operations, rigidified and pseudo-transitive, the conjugacy of the
% split tori of rank n and the problem of classifying the G between T and
% Baut(T) by an f : G --> T with three axioms. The second run runs the root
% datum out of that f: the pair (x_alpha, rho_alpha) with <rho_alpha, alpha> = 1,
% the decomposition of Lie G, the Levi decomposition (an arrow in the margin
% sends it to SGA 3 XXVI), the commutation formulas, the Weyl group W*, and on
% the last page the axioms 1) to 5) of a « système d'Enriques ».
%
% The hand is drafting, not copying out. The formulas, the underlined
% statements and the displayed conditions carry; a good deal of the connective
% prose does not, and is marked rather than filled. Three particular readings
% are flagged and not settled: the parenthesis « (bel/S) » of page 1, an
% abbreviation not resolved here; the word for the subgroups of maximal rank at
% the foot of page 1; and the two words after « X » in the phrase fixing the
% data on page 2. Page 5 carries its axioms in the upper third and is otherwise
% blank below three ruled lines.
%
% Pass: Opus 5 (claude-opus-5), 2026-09-12 — first pass, unchecked against the pages by a human.
\documentclass[11pt,a4paper]{article}
% Le préambule commun à toutes les transcriptions du fonds Grothendieck.
%
% Il tient en une idée : **l'incertitude se compose, elle ne se résout pas.**
% Une transcription qui lisse un mot douteux a détruit la seule chose pour
% laquelle on la faisait — un lecteur doit pouvoir savoir, à chaque endroit,
% ce qui était sur le feuillet et ce qui a été ajouté. Les six macros qui
% suivent sont donc l'appareil critique, et non de la décoration.
%
% Les mêmes commandes sont reconnues par `scripts/render.mjs`, qui produit la
% vue de lecture du site. Une macro ajoutée ici doit l'être aussi là-bas,
% faute de quoi le rendu échouera bruyamment — ce qui est voulu : un rendu
% silencieusement faux est pire qu'un rendu absent.

\ProvidesPackage{grothendieck}[2026/08/08 Transcriptions du fonds Grothendieck]

\usepackage{amsmath,amssymb,amsthm}
\usepackage{mathrsfs}
\usepackage{stmaryrd}
% Les diagrammes commutatifs sont partout dans ce fonds : c'est la langue
% même de la géométrie algébrique de Grothendieck, et une transcription qui
% les rendrait en prose ne transcrirait rien.
\usepackage{tikz-cd}
% Français par défaut — c'est la langue des feuillets — mais l'anglais est
% chargé aussi : la traduction et le résumé sont des documents anglais, et la
% typographie française (espaces avant les deux-points) y serait une faute.
% Ces documents basculent par \selectlanguage{english} en tête de corps.
\usepackage[english,french]{babel}
\usepackage{fontspec}
\usepackage{geometry}
\geometry{a4paper,margin=25mm,marginparwidth=28mm}
\usepackage{marginnote}
\usepackage{xcolor}
% ulem, not soul. `soul`'s \st and \ul are built on a character-by-character
% scan that XeLaTeX's font machinery breaks: the document fails with an
% undefined control sequence rather than degrading. ulem does the same two jobs
% and is Unicode-engine safe. `normalem` stops it hijacking \emph, which the
% transcriptions use for Grothendieck's own emphasis.
\usepackage[normalem]{ulem}
\usepackage{hyperref}
\hypersetup{colorlinks=true,linkcolor=black,urlcolor=black,pdfborder={0 0 0}}

% --- Métadonnées du lot -----------------------------------------------------
% Reprises telles quelles par le script de rendu, qui les affiche en tête de
% la vue de lecture. Elles fixent ce que le fichier prétend couvrir : sans
% elles, une transcription flotte sans point d'attache dans un fonds de seize
% mille feuillets.

\newcommand{\folder}[1]{\gdef\grothFolder{#1}}
\newcommand{\batch}[1]{\gdef\grothBatch{#1}}
\newcommand{\leaves}[2]{\gdef\grothFirst{#1}\gdef\grothLast{#2}}
\newcommand{\foldertitle}[1]{\gdef\grothFoldertitle{#1}}
\gdef\grothFolder{?}\gdef\grothBatch{?}\gdef\grothFirst{?}\gdef\grothLast{?}\gdef\grothFoldertitle{}

% --- Le résumé --------------------------------------------------------------
%
% Il ouvre la lecture modernisée, sous le titre « Résumé », et il est composé
% en petit. Ce n'est pas un ornement : c'est le seul passage du document qui
% ne soit pas des mathématiques, et un lecteur doit pouvoir le reconnaître
% avant de l'avoir lu — puis le sauter s'il connaît déjà le sujet, ou s'y
% arrêter s'il ne le connaît pas. Le filet qui le suit marque la reprise du
% corps.

\newenvironment{resume}
  {\small\setlength{\parskip}{0.45em}}
  {\par\medskip\hrule\medskip}

% --- Mots-clés ---------------------------------------------------------------
%
% En anglais, en clôture du résumé de la lecture modernisée : le vocabulaire
% moderne sous lequel chercher ce que les feuillets construisent. Le manifeste
% du site (`npm run manifest`) les extrait pour étiqueter la cote entière —
% cette ligne est la source unique des tags, il n'y a pas de fichier à côté.
\newcommand{\keywords}[1]{\par\smallskip\noindent
  {\footnotesize\textsc{Keywords} --- \textit{#1}}\par}

% --- L'appareil critique ----------------------------------------------------

% Le numéro de feuillet, en marge. C'est l'ancre qui permet de régler un
% désaccord sur une formule en regardant *un* feuillet plutôt qu'une chemise
% de six cents. Le site s'en sert aussi pour tourner les pages du fac-similé
% quand on fait défiler la transcription.
\newcommand{\leaf}[1]{%
  \marginnote{\footnotesize\color{gray!70}\textsf{#1}}%
  \label{leaf:#1}\ignorespaces}

% Illisible : on ne devine pas. Un mot inventé qui se lit comme les autres est
% le pire résultat possible.
\newcommand{\ill}{\textcolor{red!60!black}{[\,\ldots\,]}}

% Lecture douteuse : le mot est donné, et signalé comme incertain.
\newcommand{\uncertain}[1]{\uline{#1}}

% Ajout de l'éditeur — une lettre manquante, un mot sous-entendu.
\newcommand{\add}[1]{\textcolor{blue!50!black}{[#1]}}

% Biffé par Grothendieck lui-même. Ce qu'il a rayé fait partie du document :
% une rature dit souvent où la pensée a changé de direction.
\newcommand{\struck}[1]{\sout{#1}}

% Note du transcripteur, sur l'état matériel du feuillet.
\newcommand{\note}[1]{\par\smallskip{\small\color{gray!60!black}\textit{[#1]}}\par\smallskip}

% Note marginale de Grothendieck, distincte du corps du texte.
\newcommand{\marginal}[1]{\par\smallskip\noindent\hspace{1em}%
  {\small\color{gray!60!black}#1}\par\smallskip}

% --- Titre ------------------------------------------------------------------

\AtBeginDocument{%
  \begin{center}
    {\large\bfseries Fonds Alexandre Grothendieck --- cote n\textsuperscript{o}~\grothFolder}\\[2pt]
    {\itshape\grothFoldertitle}\\[4pt]
    {\small feuillets \grothFirst--\grothLast\ (lot \grothBatch)}\\[2pt]
    {\footnotesize\color{gray!60!black}Transcription établie d'après le fac-similé numérisé
     par l'Université de Montpellier.\\
     Les lectures incertaines sont soulignées, les illisibles notées [\,\ldots\,],
     les ajouts entre crochets.}
  \end{center}
  \vspace{1em}}

\endinput


\folder{28}
\batch{1}
\pages{1}{10}
\dating{[vers 1970]}
\watermark{Édition de démonstration}
\foldertitle{[Groupes algébriques (SGA 3)] : notes manuscrites (s.d.)}

\begin{document}

\section*{Demazure 5.1.1970}

\note{titre de sa main, souligné, en tête du premier feuillet ; à droite du
titre, isolé et souligné, « $\mathrm{Aut}^{\circ}$ ». Le titre bracketé du
fonds est celui des archivistes ; les deux dates du dossier sont les siennes}

\page{1}

\subsection*{Introduction}

\[
\begin{array}{ll}
\mathbf{P}_1 \times \mathbf{P}_1 & \mathrm{PGL}_2 \times \mathrm{PGL}_2 \\[2pt]
\mathbf{P}_2 & \mathrm{PGL}_3 \\[2pt]
F_n \quad n \geq 2 & \bigl(\mathrm{GL}(2)\cdot \Gamma(\mathcal{O}(n))\bigr)/\mu_n
\end{array}
\qquad
\begin{array}{l}
(F_0 = \mathbf{P}_1 \times \mathbf{P}_1) \\[2pt]
(F_1 = \mathbf{P}_2 \text{ éclaté})
\end{array}
\]

\note{le chiffre après $\geq$ est surchargé ; on lit $2$, lecture que les deux
parenthèses de droite appellent. Après $\mathrm{PGL}_3$ une parenthèse contient
un mot raturé. La troisième ligne de droite est écrite par-dessus une première
rédaction : le facteur $\mathrm{GL}(2)$ est repassé, et la parenthèse fermante
est doublée}

\emph{Th.} (Enriques) tout sous-groupe de $\mathrm{Cr}_2$, algébrique linéaire
connexe, est $\subset$ un des précédents : à conjugaison près \ill{}

\emph{Cor.}
\begin{itemize}
  \item \struck{\ill{}} structure réductive de type $A_n$
  \item ils sont définis sur $\mathbf{Z}$
  \item sont des groupes d'automorphismes de variétés très simples, proj. et
        \uncertain{lisses} (opérant avec orbites ouvertes)
  \item tores de rang 2
\end{itemize}

\note{le premier article porte « structure » écrit au-dessus d'un mot biffé, et
« type » est souligné ; les quatre articles sont pris dans une accolade}

On va généraliser à la dimension $n$, sauf difficultés liées aux tores
maximaux ; ou si on prend les \uncertain{s.-gpes} de rang maximum
\struck{\ill{}}

\subsection*{Sorites}

(Cas $X$ lisse \uncertain{séparé} $/S$) \note{un mot est inséré au-dessus de la
barre oblique, non lu}

Pseudo-morphismes, pseudo-isomorphismes \uncertain{(bel/$S$)}
\note{l'abréviation entre parenthèses n'est pas résolue ; les lettres se lisent
« bel », le tout suivi de « /S »}

Pseudo-automorphismes
\[
  \mathrm{Psaut}_S(X), \quad \mathrm{Baut}_S(X)
\]
\[
  \mathrm{Psaut}_k(K) ; \qquad
  \mathrm{Cr}_{n,k} = \mathrm{Baut}_k\bigl(k[t_1,\ldots,t_n]\bigr)
\]
\[
  \mathrm{Psaut}_k(K)(k) \simeq \mathrm{Aut}_k(K)^{\circ}, \qquad
  \mathrm{Lie}\bigl(\mathrm{Baut}_k(K)\bigr) \simeq \mathrm{Der}_k(K)^{\circ}
\]

\note{« Baut » est souligné à chacune de ses occurrences, « Psaut » ne l'est
pas}

\emph{Proposition} Section unité de $\mathrm{Baut}_S(X)$ immersion fermée dans
chacun des 2 cas
\begin{itemize}
  \item[a)] $S$ artinien
  \item[b)] fibres irréductibles \note{« (géom.) » ajouté au-dessous, appelé
        par un signe d'insertion}
\end{itemize}

\page{2}

\emph{Cor.} $G \longrightarrow \mathrm{Baut}_S(X)$ \note{au-dessus de la
flèche : « hom. de gp. » ; au-dessous du $G$ : « schéma en groupes »} alors le
noyau est un \ill{} \uncertain{sa réduction} \ill{} un groupe fermé

\struck{\emph{Prop.} $X$ courbe irréductible propre, $S$ loc. noeth. normal,
$\mathrm{Baut}_S(X)$ \ill{}}
\note{tout ce bloc est biffé ; à sa droite, cerné d'un trait et appelé par une
flèche, un ajout dont on lit « \uncertain{dans le corps} »}

\emph{Lemme} $G \longrightarrow \mathrm{Baut}_k(X)$, $X$ courbe lisse propre
\note{au-dessous du $G$ : « schéma en groupes »} sur $k$, il se factorise par
\struck{$\mathrm{Baut}$} $\mathrm{Aut}_k(X)$.

\ill{} \uncertain{bien} corps $k$, et $X$ \ill{} \ill{} t.f. $(/k)$. Que
signifie
\[
  G \longrightarrow \mathrm{Psaut}_k(X) \; ?
\]
\note{au-dessous de la flèche : « schéma en groupes lisse de t.f. ». Les deux
mots qui suivent le $X$ ne sont pas lus}

\[
  G \times X \dashrightarrow X
\]
ouvert de définition $U$ (un ouvert de $G \times X$, \ill{} en général
\ill{} $e \times X$)

(PO 1) \uncertain{pr. et op.} dominants

(PO 2) associativité

\begin{tikzcd}[column sep=large]
G \times G \times X \arrow[r, "u_a \times X"] \arrow[d, "G \times \varphi"'] & G \times X \arrow[d, "\varphi"] \\
G \times X \arrow[r, dashed] & X
\end{tikzcd}

\note{l'étiquette de la flèche du haut est de lecture douteuse : on lit
« $u_a \times X$ »}

\[
  \left.
  \begin{array}{l}
    g, h \in G(S) \\
    x \in X(S)
  \end{array}
  \right\}
  \quad
  \begin{array}{l}
    h \cdot x \text{ défini} \\
    g(h \cdot x) \text{ défini}
  \end{array}
  \;\Longrightarrow\;
  \bigl(gh \cdot x \text{ défini et égal à } g(h\,x)\bigr)
\]

$\exists\, N \subset G$ ($= \ker u$) tel que \ill{} $G/N \hookrightarrow
\mathrm{Baut}_k(X)$.
\note{un mot inséré au-dessus d'un mot biffé, à la fin de l'avant-dernière
ligne, n'est pas lu}

\page{3}

\begin{tikzcd}
G \times X \arrow[r, dashed] & X \times X \\
U \arrow[u, hook] & \\
F \arrow[u, hook] \arrow[r] & X \arrow[uu]
\end{tikzcd}

Opération \emph{rigidifiée} si $F \to X$ fid. plat

\emph{Prop.} Si l'opération est rigidifiée, $F_X$ est un sous-schéma en groupes
de $G_X$.

\emph{Cor.} $N =$ le plus grand sous-schéma en groupes constant de $F_X$.

Pseudo-opération (pseudo-)\emph{transitive} si $G \times X \to X \times X$
dominant

(dans le cas d'une vraie opération, cela signifie qu'il y a une (unique) orbite
dense)

\emph{Proposition} Si $G$ est pseudo-transitif sur $X$, $X$ est
pseudo-isomorphe à un espace homogène.

\subsection*{Sous-tores de $\mathrm{Cr}_{n,k}$}

\emph{Prop.} Soit $T$ tore (\uncertain{déployé}) qui pseudo-opère de façon
\struck{\ill{}} rigidifiée, ($X$ irréductible). Alors $F = e \times X$
\struck{\ill{}}

\emph{Cor. 1} $\dim T \leq \dim X$, égalité $\Longleftrightarrow$
pseudo-transitivité \note{sous « égalité » : « $T$ déployé » ; sous
« pseudo-transitivité », souligné : \uncertain{torseurs sous}}

\emph{Cor. 2} Supposons $\dim T = \dim X$, alors $X \simeq T$ (à
pseudo-isomorphisme près)

\emph{Cor. 3} $X$ irréductible, \ill{} conditions équivalentes
\begin{itemize}
  \item[a)] $X$ rationnel
  \item[b)] $\mathrm{Psaut}(X) \simeq \mathrm{Cr}_{n,k}$
  \item[c)] $\mathrm{Psaut}(X) \supset T$ tore de dim $n$
\end{itemize}

\page{4}

\emph{Cor. 4} Les tores déployés \add{de dim $n$} de $\mathrm{Cr}_{n,k}$ sont
conjugués.

\emph{Cor. 5} Ils sont maximaux, ils sont leurs propres \ill{}, et le
normalisateur est un groupe \struck{\ill{}} lisse, et $N/T \simeq
\mathrm{Aut}_{\mathrm{gp}}(T)$.

\emph{Cor. 6} $T$ est un tore maximal de $\mathrm{Cr}_{n,k}$

\emph{Proposition} $T$ déployé, $X$ irréductible, pseudo-opération pseudo
fidèle. Alors $X \dashrightarrow T \times Y$, où $Y$ est \ill{} irréductible
(quelconque) : pseudo-isomorphisme.

\emph{Cor. 1} Soient $0 \leq d \leq n$. Conditions équivalentes
\begin{itemize}
  \item[a)] $\mathrm{Im}\ \mathbf{G}_m^d \subset \mathrm{Cr}_{n,k}$ \ill{}
        contenu dans un \ill{}
  \item[b)] \ill{} extension $K/k$ de deg. tr. $d$, et $K[t_1,\ldots,t_{n-d}]$
        pure, \ill{} pure (« Lüroth »).
\end{itemize}

\begin{quote}
[NB \quad OK si $d = 0$, $d = 1$, $d = 2$ \note{en regard : « car. nulle »}
(Castelnuovo). $d = n$ \note{en regard : « $k$ alg. clos »}. Marche si
$n \leq 3$]
\end{quote}

\emph{Pb.} Classifier les $G$ alg. lisses $\subset \mathrm{Cr}_{n,k}$ qui
contiennent un tore déployé de dimension $n$.

\emph{Th. de structure} On se fixe $T$, et
\[
  T \subset G \subset \mathrm{Baut}(T)
\]
L'opération de $G$ est transitive. On veut exprimer \ill{} homogène $X$ sous
$G$, de dimension $n$, de telle façon que $T$ opère fidèlement avec une orbite
ouverte

\page{5}

Utilisant \ill{}, on aura
\[
  \begin{array}{c}
    T \subset G \\
    \cup \\
    H
  \end{array}
  \qquad
  \left\{
  \begin{array}{l}
    T \times H \longrightarrow G \quad \text{immersion ouverte} \\
    \struck{\ill{}} \quad \bigcap \text{ conjugués de } H = e
  \end{array}
  \right.
\]

On aura pseudo-explicitement
\[
  G \dashrightarrow T, \qquad (g, t) \mapsto f(gt)
\]
\note{le $f$ est porté au-dessus de la flèche pointillée ;
au-dessous de cette ligne, sans liaison explicite : « $t\,f(t^{-1}\,g\,t)$ »}

[pseudo-opération \ill{}]

\subsection*{Axiomes sur $f$}

\begin{enumerate}
  \item[1)] $f(gg') = f\bigl(g\,f(g')\bigr)$ \quad (si $f(g')$ est défini, les
        2 sens le sont simultanément)
  \item[2)] $f$ définie en tt pt de $T$, et y induit l'identité
  \item[3)] Si $g, g' \in G(S)$, et $f(gt) = f(g't)$ \ill{}, alors $g = g'$.
\end{enumerate}
\[
  \Longrightarrow \quad f(tg) = t\,f(g)
\]

\note{le reste du feuillet est blanc, sous trois longs traits horizontaux}

\section*{Demazure 12.2.70}

\page{6}

\[
  \begin{array}{c}
    T \subset G \\
    \subset G'
  \end{array}
\]
\note{au-dessus du premier $\subset$, une flèche rentrante étiquetée $f$}
possibilité de restreindre à un $G'$

\emph{Ex.} $G' = \mathrm{Cent}_G(Q)$ \qquad $Q =$ sous-tore de $T$

\note{à la ligne suivante, « N.B. » suivi de deux mots non lus et de
« central »}

Si $Z \subset T$ central dans $G$, alors \ill{} par passage au quotient, les
conditions 1) 2) 3) restent vérifiées (il faut un peu d'attention pour
$3^{\circ}$)

\subsection*{Cas particulier}

$\dim G - \dim T = 1$. $G$ n'est pas commutatif ; produit $\frac{1}{2}$ direct
donné par \ill{} non trivial (car $T$ n'est pas central si $T \neq G$).
\note{« $\frac{1}{2}$ direct » : semi-direct, son abréviation}

On aura $G = T \cdot G_u$, $G_u \simeq \mathbf{G}_a$

\emph{Proposition}
\[
  \exists\,!
  \left\{
  \begin{array}{l}
    x_\alpha : \mathbf{G}_a \xrightarrow{\ \sim\ } U \\
    \rho_\alpha : \mathbf{G}_m \longrightarrow T
  \end{array}
  \right.
\]
\note{une flèche courbe part de $T$ et revient sous les deux lignes, étiquetée
$\alpha$}
t.q.
\[
  \left\{
  \begin{array}{l}
    t\,x_\alpha(\lambda)\,t^{-1} = x_\alpha\bigl(\alpha(t)\lambda\bigr) \\
    f\bigl(x_\alpha(\lambda)\bigr) = \rho_\alpha(1+\lambda)
  \end{array}
  \right.
\]
($f$ définie \ill{} $x_\alpha(\lambda)$ ssi $1+\lambda$ inversible)

De plus, on aura
\[
  \left\{
  \begin{array}{l}
    \langle \rho_\alpha, \alpha \rangle = 1 \\
    H = x_\alpha(1)\,T\,x_\alpha(-1) \qquad (?)
  \end{array}
  \right.
\]
\note{le « (?) » est de sa main}

Inversement, si on se donne $\alpha, \rho_\alpha$ tels que $\langle \alpha,
\rho_\alpha \rangle = 1$, on reconstitue \ill{}

\subsection*{Cas général}

$T \subset G$ \note{$f$ porté au-dessous, en pointillé}
\[
  \mathfrak{g} = \mathfrak{g}_0 + \sum_{\alpha \in R} \mathfrak{g}^{\alpha},
  \qquad R \subset \mathrm{Hom}(T, \mathbf{G}_m)
\]
\[
  \mathfrak{g}_0 = \mathrm{Lie}\,T \quad (T = \mathrm{Cent}_G T), \qquad
  \dim \mathfrak{g}^{\alpha} = 1
\]

\page{7}

Soit $\chi : T \to \mathbf{G}_m$
\[
  T_\chi = (\mathrm{Ker}\,\chi)^{\circ} \subset T, \qquad
  Z_\chi = \mathrm{Cent}(T_\chi)
\]
$Z_\chi / T_\chi$ a comme tore maximal $T/T_\chi$ qui \ill{} de dim 1
\[
  Z_\chi / T_\chi \hookrightarrow \mathrm{Aut}(\mathbf{P}^1) = \mathrm{PGL}(2)
  \qquad (\text{car } \mathbf{P}^1 \text{ de dim } 1)
\]
\note{une flèche montante joint cette ligne à « $\mathrm{Baut}(G/T_\chi)$ »,
porté au-dessus}

4 possibilités, \uncertain{suivant le type} de $\mathrm{PGL}(2)$ \ill{} tore
maximal :
\begin{itemize}
  \item tore maximal \quad : \quad $Z_\chi = T$
  \item \ill{} non-commutatives \quad : \quad $Z_\chi/T$ de rang 2
  \item Borel \quad : \quad $\dim Z_\chi/T = 1$
  \item tout \quad : \quad $Z_\chi$ \uncertain{identifié de rg.} \ill{}
\end{itemize}
\note{en regard du dernier article, encadré : « $\frac{1}{2}$ \ill{} 1 »}

\emph{Conclusion} Supposons $\chi \in R$, les deux premières sont exclues.

\begin{enumerate}
  \item[(1)] On a
  \[
    \left\{
    \begin{array}{l}
      \dim \mathfrak{g}_\alpha = 1 \\
      \exists\,! \; x_\alpha : \mathbf{G}_a \longrightarrow Z_\chi \subset G,
        \quad x_\alpha(\mathbf{G}_a) = U_\alpha, \;
        \mathrm{Lie}(U_\alpha) = \mathfrak{g}_\alpha \\
      \phantom{\exists\,! \;} t\,x_\alpha(\lambda)\,t^{-1} =
        x_\alpha\bigl(\alpha(t)\lambda\bigr) \\
      \rho_\alpha : \mathbf{G}_m \longrightarrow T, \quad
        f\bigl(x_\alpha(\lambda)\bigr) = \rho_\alpha(1+\lambda), \quad
        \langle \rho_\alpha, \alpha \rangle = 1
    \end{array}
    \right.
  \]
  \[
    \mathrm{Lie}\,Z_\chi = \mathfrak{g}_0 +
      \sum_{\beta \in \mathbf{Q}\chi \,\cap\, R} \mathfrak{g}^{\beta}
  \]
  \item[(2)] Si $\lambda\alpha \in R$, $\lambda \neq 1$, alors
    $\lambda = -1$ et $Z_\alpha$ est \uncertain{identifié de} \uncertain{19.SS.1},
    d'où les 2 sous-groupes de Borel qui $\supset T$, à savoir $T\,U_\alpha$ et
    $T\,U_{-\alpha}$
\end{enumerate}
\note{les deux numéros sont cerclés de sa main ; la référence donnée pour
$Z_\alpha$ est de lecture douteuse}

\note{une double barre suivie d'une flèche descendante porte, écrit dessus :
« SGA 3 \uncertain{XXVI} », doublement souligné}

$\exists$ décomposition de Lévi
\[
  G^{\circ} = L \cdot U, \qquad L \supset T
\]
\note{sous $L$, appelé par une flèche : « réductif »}
Les racines de $L$ sont les $\alpha \in R$ telles que $-\alpha \in R$.

\page{8}

\emph{Cor.} \add{si $G$ connexe,} $G^{\circ}$ est engendré par les $U_\alpha$
et $T$ (\uncertain{groupe abélien}) ; $f$ est déterminée par la connaissance
des $x_\alpha$, $\rho_\alpha$

\emph{Th.} Si $\alpha, -\alpha \in R$, $\exists\, \varphi_\alpha :
\mathrm{GL}(2) \longrightarrow G$ tel que
\[
  \varphi_\alpha \begin{pmatrix} 1 & \lambda \\ 0 & 1 \end{pmatrix}
    = x_\alpha(\lambda), \qquad
  \varphi_\alpha \begin{pmatrix} 1 & 0 \\ \lambda & 1 \end{pmatrix}
    = x_{-\alpha}(\lambda)
\]
\[
  \varphi_\alpha \begin{pmatrix} \mu & 0 \\ 0 & 1 \end{pmatrix}
    = \rho_\alpha(\mu), \qquad
  \varphi_\alpha \begin{pmatrix} 1 & 0 \\ 0 & \mu \end{pmatrix}
    = \rho_{-\alpha}(\mu)
\]
\[
  f\left(\varphi_\alpha \begin{pmatrix} a & b \\ c & d \end{pmatrix}\right)
    = \rho_\alpha(a+b)\,\rho_{-\alpha}(c+d)
\]
\[
  \varphi_\alpha \begin{pmatrix} a & b \\ c & d \end{pmatrix} \cdot t
    = t \cdot \rho_\alpha\!\left(a + \frac{b}{\alpha(t)}\right)
      \rho_{-\alpha}\bigl(c\,\alpha(t) + d\bigr)
\]

\uncertain{Obs.} Considérons $\varphi_\alpha \begin{pmatrix} 0 & 1 \\ 1 & 0
\end{pmatrix} = w_\alpha$, on trouve
\[
  w_\alpha \cdot t = \struck{\ill{}} \; s_\alpha(t)
    = t\;\alpha^{*}\bigl(\alpha(t)\bigr)^{-1}, \qquad
  \alpha^{*} = \rho_\alpha - \rho_{-\alpha}
\]
\note{sous $\alpha^{*}$, appelé par un trait : « coracine »}

\[
  \left\{
  \begin{array}{l}
    \text{Si } M = \mathrm{Hom}(T, \mathbf{G}_m), \; \ill{} \text{ syst. de
      racines de } L \text{ est } \{M, R_s, \; \alpha \mapsto \rho_\alpha -
      \rho_{-\alpha}\} \\
    \phantom{\text{Si } M} R_s = R \cap (-R)
  \end{array}
  \right.
\]
\[
  \left\{
  \begin{array}{l}
    \mathrm{Norm}_{G^{\circ}}(T) = T \cdot W^{*} \\
    W^{*} \text{ engendré par les } w_\alpha \; (\alpha \in R_s), \quad
      R_s = R \cap (-R)
  \end{array}
  \right.
\]
\note{« $R \cap (-R)$ » est écrit au-dessous de $R_s$, relié par un signe
d'égalité, aux deux endroits}

\note{en marge droite : « NB Ce $W^{*}$ opère trivialement \ill{} sur $T$ ! »}

$W^{*}$ opère sur $T$ \ill{} \struck{\ill{}} $R$ et on aura
\[
  \rho_{s_\alpha(\beta)} = s_\alpha(\rho_\beta), \qquad
  \mathrm{int}(w_\alpha)\,x_\beta = x_{s_\alpha(\beta)}
\]

\page{9}

\emph{Proposition} Soient $\alpha, \beta \in R$ ; $U_\alpha$ et $U_\beta$
commutent \ill{}
\[
  \rho_\alpha = \rho_\beta \quad \text{ou} \quad
  \langle \rho_\alpha, \beta \rangle = \langle \rho_\beta, \alpha \rangle = 0
\]

\emph{Cor 1} $\rho_\alpha \neq \rho_\beta \Longrightarrow \langle \rho_\alpha,
\beta \rangle \leq 0$

\emph{Cor. 2} $\langle \rho_\alpha, \beta \rangle < 0$, $\langle \rho_\beta,
\alpha \rangle < 0 \Longrightarrow \alpha + \beta = 0$

\struck{Si $\rho_\alpha \neq \rho_\alpha$}

\emph{Cas}
\begin{itemize}
  \item[a)] $\rho_\alpha = \rho_\beta$
  \item[b)] $\alpha + \beta = 0$
  \item[c)] à l'index près, \quad $\left\{
    \begin{array}{l}
      \langle \rho_\alpha, \beta \rangle = -n \leq 0 \\
      \langle \rho_\beta, \alpha \rangle = 0
    \end{array} \right.$
\end{itemize}

\emph{Prop.} Dans le cas c) ci-dessus :
\begin{enumerate}
  \item[1)] $\lambda\alpha + \mu\beta \subset R$, $\lambda, \mu \in \mathbf{N}
    \Longrightarrow \lambda\alpha + \mu\beta = \left\{
    \begin{array}{l} \alpha \\ \beta + i\alpha \end{array} \right.$
    \; où $i \leq$ \ill{}
  \item[2)] Si $\exists\, 0 \leq i \leq n$, et si $\binom{n}{i} \neq 0$ dans
    $k$, alors $\beta + i\alpha \in R$ \note{au-dessus : « et
    $\rho_{\beta+i\alpha} = \rho_\beta$ »}
  \item[3)] $x_\alpha(\lambda)\,x_\beta(\mu)\,x_\alpha(-\lambda) = \prod_i
    x_{\beta + i\alpha}\left(\pm \binom{n}{i} \lambda^{i} \mu\right)$
\end{enumerate}

\[
  [X_\alpha, X_\beta] = \bigl(\langle \rho_\alpha, \beta \rangle -
    \langle \rho_\beta, \alpha \rangle\bigr) X_{\alpha+\beta}
  \qquad \text{si } \alpha + \beta \neq 0
\]
de plus
\[
  [X_\alpha, X_{-\alpha}] = \rho_\alpha - \rho_{-\alpha}
\]

\emph{NB} Toutes les formules à coeff. $\in \mathbf{Z}$.

\emph{Th. de conjugaison} Deux sous-groupes $G$, $G' \supset T$ de
$\mathrm{Cr}_{n,k}$, \struck{\ill{}} $G$, $G'$ \uncertain{connexes}, avec les
mêmes \struck{\ill{}} $R$ et $\rho$, sont conjugués dans $\mathrm{Cr}_{n,k}$
\ill{}

\page{10}

\subsection*{Syst. d'Enriques}

$M$, $M^{*}$, $R \subset M$, $\rho : R \longrightarrow M^{*}$

\begin{enumerate}
  \item[1)] $R$ fini
  \item[2)] $\langle \rho_\alpha, \alpha \rangle = 1$ \qquad
    $\forall\, \alpha \in R$
  \item[3)] $\rho_\alpha \neq \rho_\beta \Longrightarrow
    \langle \rho_\alpha, \beta \rangle \leq 0$
  \item[4)] $\langle \rho_\alpha, \beta \rangle < 0$, $\langle \rho_\beta,
    \alpha \rangle < 0 \Longrightarrow \alpha + \beta = 0$
  \item[5)] Si $\rho_\alpha \neq \rho_\beta$, $\langle \rho_\beta, \alpha
    \rangle = 0$, alors $\beta + i\alpha \in R$ dès que $\binom{n}{i} \neq 0$
    dans $k$ \quad ($0 \leq i \leq n$)
\end{enumerate}
\note{un gribouillis biffé précède le numéro 5)}

\emph{Th. d'existence} On trouve tous les syst. d'Enriques.

\end{document}
